\capitulosinnumero{Resumen}

Este documento recoge la memoria del Trabajo de Fin de Grado del Grado en Ingeniería Multimedia de Leigh Garrett. El proyecto, titulado «Desarrollo de un videojuego narrativo en 2D con mecánicas adaptativas sobre la perspectiva de género» y realizado durante el curso académico 2026-2027, ha sido tutorizado por Carlos J. Villagrá Arnedo y Patricia María Compañ Rosique.

En los últimos años, los videojuegos han demostrado ser un medio capaz de contar historias muy personales y tratar temas humanos complejos. Este trabajo consiste en la idea, el diseño y el desarrollo de un prototipo de videojuego narrativo en \ac{2D}, visto en primera persona con efecto de profundidad (2.5D), que trata sobre el agotamiento mental que provoca una crisis de identidad de género debido al rechazo social e internalizado. Para transmitir esta vivencia desde la perspectiva de género, el juego comienza engañando al jugador con la premisa de un apocalipsis zombi de supervivencia, haciéndole pasar por tres mundos que van cambiando visual y mecánicamente hasta revelar que todo forma parte de la presión cotidiana en su propio instituto.

La aportación principal del proyecto a nivel técnico y de diseño es un sistema de jugabilidad adaptativa que funciona en segundo plano. A través de un módulo llamado Director de Alucinación, el juego observa cómo actúa el jugador durante la partida y registra sus decisiones en tres ejes o dimensiones continuas (Ímpetu, Cautela y Templanza), evitando así encasillarlo en perfiles cerrados o estereotipados. Con estos datos, el sistema adapta la dificultad de los escenarios, las habilidades disponibles y el aspecto de la propia interfaz (\ac{UI}) mediante «Mecanismos de Defensa» programados con \textit{ScriptableObjects}. Además, para no frustrar al jugador ni quitarle el control de forma injusta, el juego incorpora un sistema de consentimiento explícito que le pregunta antes de aplicar cualquier cambio en sus mecánicas, contando también con el asesoramiento de una especialista en psicología para tratar el tema con el cuidado necesario.

El propósito final de la partida es avanzar por este recorrido interno hasta llegar a la azotea del instituto, donde tiene lugar un desenlace abierto que varía según cómo haya afrontado el jugador cada etapa del camino. Esta memoria documenta todo el proceso previo y de diseño: desde el estudio del estado del arte y la justificación de las herramientas elegidas, hasta la elaboración del Documento de Diseño de Juego (\ac{GDD}), el apartado artístico en primera persona y la estructura técnica planteada en el motor Unity.

\vspace{1em}
\noindent\textbf{Palabras clave:} Videojuego 2D, Primera Persona 2.5D, Narrativa Interactiva, Mecánicas Adaptativas, Perspectiva de Género, Salud Mental, Identidad de Género, Unity.